\documentclass{article}

% ============================================================
% MA 213 Lab Activity Template
% Student-facing worksheet template for lab development.
% Copy this file into a lab directory, then replace bracketed
% placeholders such as [WEEK], [Lab title], and [prompt].
% ============================================================

\usepackage[margin=1in]{geometry}
\usepackage{xcolor}
\usepackage{parskip}
\usepackage{array}
\usepackage{enumitem}
\usepackage{listings}

\definecolor{codegray}{gray}{0.97}
\definecolor{huddleblue}{RGB}{214,234,255}
\definecolor{predictionyellow}{RGB}{255,242,200}
\definecolor{debatepink}{RGB}{255,225,225}

\lstset{
  basicstyle=\ttfamily\small,
  columns=fullflexible,
  frame=single,
  framerule=0.4pt,
  backgroundcolor=\color{codegray},
  breaklines=true,
  showstringspaces=false,
  keepspaces=true,
  xleftmargin=0pt,
  xrightmargin=0pt
}

\newcommand{\coursenumber}{MA 213}
\newcommand{\weekplaceholder}{[WEEK]}
\newcommand{\labplaceholder}{[LAB NUMBER]}
\newcommand{\labtitle}{[Lab title]}

\newcommand{\nameblank}{\rule{1.75in}{0.4pt}}
\newcommand{\idblank}{\rule{1.55in}{0.4pt}}
\newcommand{\shortblank}{\rule{1.6in}{0.4pt}}
\newcommand{\longblank}{\rule{0.93\linewidth}{0.4pt}}

\newcommand{\responsebox}[2]{
  \noindent\fbox{%
    \begin{minipage}{0.95\linewidth}
      \textbf{#1} #2
      \vspace{0.55in}
    \end{minipage}
  }
}

\newcommand{\linedbox}[1]{
  \noindent\fbox{%
    \begin{minipage}{0.95\linewidth}
      #1
      \vspace{0.18in}

      \longblank

      \vspace{0.18in}
      \longblank

      \vspace{0.18in}
      \longblank
    \end{minipage}
  }
}

\newcommand{\callout}[3]{
  \noindent\colorbox{#1}{%
    \makebox[\dimexpr0.95\linewidth-2\fboxsep\relax][l]{\textbf{#2}}%
  }

  \noindent\fbox{%
    \begin{minipage}{0.93\linewidth}
      #3
    \end{minipage}
  }
}

\begin{document}

\begin{center}
  {\LARGE \coursenumber{} Week \weekplaceholder{}: Lab \labplaceholder{}}\\
  {\large \labtitle{}}
\end{center}

\begin{enumerate}[label=\arabic*.,leftmargin=*,itemsep=.7em]
  \item \textbf{Your name:} \nameblank \hfill \textbf{BU ID:} \idblank
  \item \textbf{Partner's name:} \nameblank \hfill \textbf{BU ID:} \idblank
\end{enumerate}

\textbf{Big idea:} [One or two sentences explaining the statistical purpose of the lab.]
\textbf{Do not use GenAI tools for this lab.} The point is to practice your own analysis work, coding skills, and statistical reasoning.

\textbf{Files used:} \texttt{[data-file-1.csv]} [and \texttt{[data-file-2.csv]}]

\textbf{Skills practiced:} \texttt{[function()]}, \texttt{[function()]}, \texttt{[concept]}, and \texttt{[concept]}.

\textbf{Your mission:} [A student-facing description of what partners will figure out by the end of lab.]

\section*{Icebreaker}

Before opening R: introduce yourself to your partner and answer the warm-up question below.

\responsebox{Warm-up response:}{[prompt]}

\section*{Getting Started}

[Briefly tell students what to open, download, import, or discuss first.]

\begin{lstlisting}
# Load packages if needed.
library(dplyr)

# Read in the data.
data_name <- read.csv("[data-file].csv")

# First look.
head(data_name)
str(data_name)
\end{lstlisting}

\callout{huddleblue}{HUDDLE UP}{Before calculating anything, what do you think one row represents? What variables look most important for today's question?}

\newpage

\section*{Question 1. [Explore the Data]}

[State the task in plain language. Start with interpretation before calculations when possible.]

\begin{lstlisting}
# Starter code or hints go here.
dim(data_name)
names(data_name)
\end{lstlisting}

\begin{center}
\renewcommand{\arraystretch}{2.3}
\begin{tabular}{|p{0.24\linewidth}|p{0.27\linewidth}|p{0.36\linewidth}|}
\hline
\textbf{Variable or data set} & \textbf{Type or role} & \textbf{What it means} \\
\hline
\texttt{[variable 1]} & & \\
\hline
\texttt{[variable 2]} & & \\
\hline
\texttt{[variable 3]} & & \\
\hline
\end{tabular}
\end{center}

\linedbox{\textbf{Write 2-3 sentences describing what you learned from the first look at the data.}}

\section*{Question 2. [Make a Prediction]}

[Ask students to make a prediction, choose a variable, choose a year/group, or state an expectation.]

\callout{predictionyellow}{PREDICTION TIME}{Before running the code, predict what pattern you expect to see. Be specific enough that the data can support or challenge your answer.}

\begin{lstlisting}
filtered_data <- data_name %>%
  filter(________ == ________)

head(filtered_data)
\end{lstlisting}

\newpage

\linedbox{\textbf{What did your team predict? What evidence did you check first?}}

\section*{Question 3. [Compute a Summary or Create a Table]}

[Use this section for numerical summaries, grouped summaries, or contingency tables.]

\begin{lstlisting}
summary_table <- data_name %>%
  group_by(________) %>%
  summarize(
    count = n(),
    mean_value = mean(________)
  )

summary_table
\end{lstlisting}

\callout{huddleblue}{HUDDLE UP}{Does the summary answer the question directly, or do you need another comparison, plot, or table?}

\callout{huddleblue}{CLAIM, EVIDENCE, CONTEXT}{%
The main pattern is \shortblank.\\[.6em]
My evidence is \shortblank.\\[.6em]
This matters because \shortblank.
}

\section*{Question 4. [Interpret in Context]}

[Ask students to connect the output to the real-world context or course concept.]

\begin{center}
\renewcommand{\arraystretch}{2.25}
\begin{tabular}{|p{0.42\linewidth}|p{0.50\linewidth}|}
\hline
\textbf{Question} & \textbf{Your answer} \\
\hline
What comparison did you make? & \\
\hline
What result is strongest? & \\
\hline
What limitation or caution should readers know? & \\
\hline
\end{tabular}
\end{center}

\newpage

\callout{debatepink}{DEBATE CORNER}{If you and your partner disagree about the interpretation, write both arguments first. Then decide what claim the evidence actually supports.}

\section*{Question 5. [Close the Lab]}

[End with a synthesis question, short deliverable, or decision students can submit.]

\linedbox{\textbf{Final response:} state your claim, cite one piece of evidence, and explain the context.}

\section*{Optional Challenge}

[Add one extension for groups that finish early. Keep it useful but nonessential.]

\begin{lstlisting}
# Optional challenge starter code.
\end{lstlisting}

\end{document}
